\begin{answers}{ch:eval-serving}{評価と推論：モデルを測り、届ける}
\item 対応のない計算では差は標準誤差の約 1.73 倍、対応のある計算では約 3.19 倍で、結論は「差があるとは言えない」から「偶然とは考えにくい」に変わります。
正答率は C が 82.0\%、D が 86.0\% で、差は 4.0 ポイントです。
標準誤差は C が $\sqrt{0.82\times0.18/500}$ で 1.718 ポイント、D が 1.552 ポイントなので、式\eqref{eq:es-unpaired}の差の標準誤差は $\sqrt{1.718^2+1.552^2}\approx2.315$ ポイントです。
差はその 1.73 倍で、1.96 倍に届きません。
式\eqref{eq:es-paired}では $\frac{1}{500}\sqrt{40-20^2/500}$ で 1.252 ポイントになり、差はその 3.19 倍です。
符号検定の $p$ 値も 0.002 で、実力が同じなら数百回に 1 回しか起きない偏りです。
同じ問題で比べると、問題の難しさという共通のぶれが消え、食い違った 40 問だけで差を測れるからです。
\item 偶然とは考えにくい、と言えます。
実行すると、\file{ch19/out/compare_ch22.txt} と同じ出力になります。
1,000 問では、\file{sft.pt} だけ正解が 35 問、\file{ch22_rl.pt} だけ正解が 56 問です。
差の 2.1 ポイントは対応ありの標準誤差 0.95 ポイントの 2.21 倍で、符号検定の $p$ 値は 0.035 です。
1 万問でも 413 問対 519 問で、差の 1.06 ポイントは標準誤差 0.31 ポイントの 3.47 倍、$p$ 値は 0.0006 でした。
対応のない計算では、1,000 問で 1.27 倍にとどまり、見分けられません。
ただし、退けたのは「問題の選び方による偶然」だけです。
学習をやり直すだけでも 0.7 ポイントほど動くので（\secref{sec:es-compare}）、伸びのすべてが強化学習の効果だとは、まだ言えません。
\item 効果は 2 人ではまだわずかで、3〜4 人からはっきり現れます（以下は実測の一例です）。
\file{load_test.py} はサーバを自分で起動するので、\code{for slots in [1, 16]} を \code{[1, 4, 16]} に書き換え、保存先も変えて実行します。
同時 1 人では、スロットが何個でもスループットもレイテンシも変わりません。
まとめる相手がいないからです。
2 人では、4 個と 16 個のスループットは 1 個より 1 割ほど高いだけでした。
3 人で 1.3〜1.4 倍、4 人で約 1.5 倍になります。
4 個と 16 個の差は 8 人から現れ、16 人では 16 個が 1 個の 3 倍前後でした。
スロットを増やす効果が出るのは、同時に使う人数が、増やす前のスロットの数を超えてからです。
2 人で伸びが小さいのは、$B=2$ のスループットが $B=1$ の 1.3 倍しかないからです（\figref{fig:es-batch-size}）。
\item 長いプロンプトを小分けにして、ほかの列のステップの合間に少しずつ読ませる方法があります。
本書のエンジンスレッド（\coderef{code:es-engine-loop}）は、\code{admit} のプリフィルが終わってから \code{step} に進むので、デコード中の列はプリフィル 1 回の間だけ待ちます。
短いプロンプトなら 1 ステップ分ほどですが、長いほど延びます。
4 列をデコードしている途中に 120 トークンのプロンプトを入れると、その 1 回だけ、トークン間時間がふだんの約 4 倍になりました（実測の一例。計算機が混むともっと延びます）。
プロンプトを 8 トークンずつに分け、1 ステップにつき 8 トークンだけ \code{engine.forward} で読ませると、延びは約 1.8 倍に収まり、どの列の出力も変わりません。
代わりに、延びは 15 ステップ続き、長いプロンプトの利用者が 1 個目を受け取るのも 15 ステップ後です。
待ちをなくすのではなく、薄く分け合う工夫です。
前置きが共通なら、プレフィックスキャッシュで読む量を減らすこともできます（\secref{sec:es-cache}）。
\item 質の近い二つの応答の対を約 400 件用意し、どの対も順序を入れ替えて 2 回、同じ審査員に判定させます。
位置の偏りがなければ、順序を入れ替えても同じ応答が選ばれるはずです。
信頼区間を付ける数字は二つです。
一つ目は、2 回の判定が食い違った対の割合で、式\eqref{eq:es-se}と式\eqref{eq:es-ci}を使います。
標準誤差は割合が 50\% のときに最大なので、$1.96\sqrt{0.5\times0.5/n}=0.05$ を解いた $n\approx384$ 件があれば、幅は $\pm5$ ポイント以内に収まります。
二つ目は偏りの向きです。
2 回とも先に置いたほうを選んだ対の数を $c$、2 回とも後に置いたほうを選んだ対の数を $b$ として、$(c-b)/n$ に式\eqref{eq:es-paired}の標準誤差を付け、符号検定で $p$ 値を求めます。
どの応答も先と後に 1 回ずつ置かれるので、質の差は打ち消し合い、位置の効果だけが残ります。
\end{answers}
